\documentclass{article}
\usepackage[margin=1in]{geometry}
\usepackage{amsmath, amssymb, amsthm, enumitem}

\everymath{\displaystyle}

\title{How to Think Like a Mathematician}
\author{}
\date{}

\begin{document}
\maketitle
\section*{Chapter 16: How to read a theorem}
\begin{enumerate}[label=16.\arabic*]
  \item \begin{enumerate}[label=(\roman*)]
          \item \begin{enumerate}[label=(\alph*)]
                  \item `if' part is true; `only if' part is false.
                        \setcounter{enumiii}{2}
                  \item True.
                  \item True. Proof by induction.
                  \item True.
                        \begin{proof}
                          Because $n \ge 7$,
                          \begin{align*}
                                     & \frac{8n-12}{n(n-2)(n-6)}                                  & \ge & 0            \\
                            \implies & \frac{n^2-\left(n^2-8n+12\right)}{n\left(n^2-8n+12\right)} & \ge & 0            \\
                            \implies & \frac{n}{n^2-8n+12} - \frac{1}{n}                          & \ge & 0            \\
                            \implies & \frac{n}{n^2-8n+12}                                        & \ge & \frac{1}{n}. \\
                          \end{align*}
                        \end{proof}
                  \item True by induction. In the induction step, divide the the set $S_{k+1}$
                        to two subsets. The subset without the element has $|S_k|$ elements; the
                        subset with the new element is constructed by adding the new element to each
                        element in $S_k$, as a result, it also has $|S_k|$ elements.
                \end{enumerate}
        \end{enumerate}
\end{enumerate}

\section*{Chapter 20: Techniques of proof I: Direct method}
\begin{enumerate}[label=20.\arabic*]
  \setcounter{enumi}{5}
  \item \begin{enumerate}[label=(\roman*)]
          \item We want to show that $(-x)\times(-y)>0$ if $x$ and $y$ are positive. We have $(-x)
                  \times(-y) = (-1 \times x)\times(-1 \times y) = (-1)\times(-1) \times xy = 1
                  \times xy$ (by the proven collorary). Because $x$ and $y$ are positive, $1 \times
                  xy = xy$ is clearly positive. Thus, $(-x)\times(-y)>0$.
        \end{enumerate}
        \setcounter{enumi}{13}
  \item \begin{enumerate}[label=(\roman*)]
          \setcounter{enumii}{6}
          \item Because $x<y$, $y-x>0$. Suppose that $y-x=\delta>0$. Consider the number $z=x+\frac
                  {\delta}{\delta+1}$. This number is clearly a rational number. In addition,
                because $\frac{\delta}{\delta+1}>0$, $x > z$. Also, $\frac{\delta}{\delta+1}<\delta
                  \implies z < y$. Hence, $\exists z\in\mathbb{Q}: x<z<y$.
                \setcounter{enumii}{9}
          \item The product is odd if and only if none of the factor is even. Consider the factor
                $x_i-i$. This factor is odd if and only if one term is odd and the other is even.
                Thus, we need an arrangement such that for all factors $x_i-i$, either $x_i$ is odd
                and $i$ is even or $x_i$ is even and $i$ is odd. If $n$ is odd, for $x_i$ and $i$,
                there are exactly $\frac{n-1}{2}$ evens and $\frac{n+1}{2}$ odds. Thus, in any
                arrangements of matching $x_i$ with $i$, there will be at least one $x_i-i$ whose
                both terms are odd. Thus, this $x_i-i$ factor is even. As a result, the product is
                even.
        \end{enumerate}
\end{enumerate}
\end{document}
